\section*{अध्याय \ref{ch:g5:weather-and-climate} --- मौसम और जलवायु}

\begin{solution}{exo:g5:weather-and-climate:1}
(क) \omterm{def:g1:watching-the-weather:weather}{मौसम} (केवल आज सुबह)। (ख) \omterm{def:g5:weather-and-climate:climate}{जलवायु} (हर गर्मी)। (ग) \omterm{def:g1:watching-the-weather:weather}{मौसम}
(कल दोपहर बाद)। (घ) \omterm{def:g5:weather-and-climate:climate}{जलवायु} (हर महीने, साल-दर-साल)।
\end{solution}

\begin{solution}{exo:g5:weather-and-climate:2}
किसी जगह की \omterm{def:g5:weather-and-climate:climate}{जलवायु} वह \omterm{def:g1:watching-the-weather:weather}{मौसम} है जो वहाँ आमतौर पर रहता है, महीने-दर-महीने, बहुत
वर्षों तक। मौसम वैज्ञानिक 30 वर्षों के मापों का औसत निकालते हैं।
\end{solution}

\begin{solution}{exo:g5:weather-and-climate:3}
सबसे ठंडा महीना दिसंबर है (लगभग 6~$^\circ$C; जनवरी भी लगभग उतनी ही ठंडी है)।
सबसे गीला महीना नवंबर है (लगभग 160~मिमी)। नहीं: वक्र कभी 0~$^\circ$C से नीचे
नहीं जाता।
\end{solution}

\begin{solution}{exo:g5:weather-and-climate:4}
जुलाई सबसे गरम महीना है, लगभग 5~$^\circ$C पर। कोई भी महीना
10~$^\circ$C से ज़्यादा गरम नहीं है।
\end{solution}

\begin{solution}{exo:g5:weather-and-climate:5}
\omterm{def:g3:continents-and-oceans:equator}{भूमध्य रेखा}: उष्ण कटिबंध। \omterm{def:g3:continents-and-oceans:poles}{उत्तरी ध्रुव}: शीत कटिबंध (उत्तरी वाला)। \omterm{def:g3:continents-and-oceans:equator}{भूमध्य रेखा}
और \omterm{def:g3:continents-and-oceans:poles}{दक्षिणी ध्रुव} के ठीक आधे रास्ते पर: दक्षिणी शीतोष्ण कटिबंध (वह \omterm{def:g5:weather-and-climate:tropics}{मकर रेखा} और
अंटार्कटिक वृत्त के
बीच है)।
\end{solution}

\begin{solution}{exo:g5:weather-and-climate:6}
हिलो: उष्ण कटिबंध। फ़ीनिक्स और सिएटल: उत्तरी शीतोष्ण कटिबंध।
उत्कियाग्विक: उत्तरी शीत कटिबंध, आर्कटिक वृत्त के आगे।
\end{solution}

\begin{solution}{exo:g5:weather-and-climate:7}
पूरा साल: $147 + 96 + 106 + 81 + 48 + 37 + 15 + 25 + 41 + 99 + 160 + 145
= 1\,000$ मिलीमीटर। दिसंबर, जनवरी और फ़रवरी:
$145 + 147 + 96 = 388$ मिलीमीटर, यानी तीन महीनों में साल की \omterm{def:g4:clouds-rain-snow:precipitation}{बारिश} के एक-तिहाई से
ज़्यादा।
\end{solution}

\begin{solution}{exo:g5:weather-and-climate:8}
फ़ीनिक्स: $35 - 14 = 21$ डिग्री। हिलो: $25 - 22 = 3$ डिग्री। ऋतुएँ फ़ीनिक्स में
तापमान को बहुत ज़्यादा बदलती हैं; हिलो में वह पूरे साल लगभग एक-सा
रहता है।
\end{solution}

\begin{solution}{exo:g5:weather-and-climate:9}
\omterm{def:g5:weather-and-climate:rainforest}{वर्षावन} \omterm{def:g5:weather-and-climate:climate}{जलवायु}: पूरे साल गरमी रहती है (27 या 28~$^\circ$C, बिना सर्दी के), और
हर महीने बहुत \omterm{def:g4:clouds-rain-snow:precipitation}{बारिश} गिरती है (कम से कम 160~मिमी, और साल में 2\,000~मिमी से
ज़्यादा)। यह शहर सबसे ज़्यादा संभावना से उष्ण कटिबंध में, \omterm{def:g3:continents-and-oceans:equator}{भूमध्य रेखा} के पास
है।
\end{solution}

\begin{solution}{exo:g5:weather-and-climate:10}
पहली व्याख्या: \omterm{def:g2:water-on-earth:ocean}{समुद्र} \omterm{def:g1:landscapes-around-us:water}{तट} उष्ण कटिबंध में है, जहाँ पूरे साल गरमी रहती है, और \omterm{def:g4:clouds-rain-snow:precipitation}{हिम}
से ढका बग़ीचा किसी शीतोष्ण या शीत कटिबंध में, जहाँ जनवरी सर्दी में आती है। दूसरी
व्याख्या: \omterm{def:g2:water-on-earth:ocean}{समुद्र} \omterm{def:g1:landscapes-around-us:water}{तट} पृथ्वी के दक्षिणी आधे हिस्से में है, जहाँ जनवरी गर्मी में आती
है, और \omterm{def:g4:clouds-rain-snow:precipitation}{हिम} से ढका बग़ीचा उत्तरी आधे हिस्से में, जहाँ जनवरी सर्दी में आती है: \omterm{def:g3:continents-and-oceans:globe}{ग्लोब}
के दोनों आधे हिस्सों के बीच ऋतुएँ
उलटी होती हैं।
\end{solution}

\begin{solution}{exo:g5:weather-and-climate:11}
तापमान: फ़ीनिक्स बहुत ज़्यादा गरम है, जुलाई में लगभग 35~$^\circ$C, जबकि सिएटल में
अगस्त में लगभग 20~$^\circ$C, और फ़ीनिक्स की जनवरी (लगभग
14~$^\circ$C) सिएटल के मई जितनी गरम है। \omterm{def:g4:clouds-rain-snow:precipitation}{बारिश}: फ़ीनिक्स में, जो एक \omterm{def:g5:weather-and-climate:desert}{मरुस्थल} है,
साल में 200~मिमी से कम, जबकि सिएटल में लगभग 1\,000~मिमी। इसलिए उनकी \omterm{def:g5:weather-and-climate:climate}{जलवायु} एक-सी
नहीं है। कटिबंध हमें मोटे तौर पर बताते हैं कि \omterm{def:g3:continents-and-oceans:equator}{भूमध्य रेखा} से \omterm{def:g3:continents-and-oceans:poles}{ध्रुवों} तक सूर्य किसी
जगह को कितना गरम करता है; वे यह नहीं बताते कि कितनी \omterm{def:g4:clouds-rain-snow:precipitation}{बारिश} गिरती है, न \omterm{def:g2:water-on-earth:ocean}{समुद्र},
\omterm{def:g1:landscapes-around-us:landform}{पर्वतों} या हवाओं का असर।
\end{solution}
